\documentclass[11pt,a4paper]{article}
\usepackage[margin=1in]{geometry}
\usepackage{amsmath,amssymb}
\usepackage{booktabs}
\usepackage{enumitem}
\usepackage{hyperref}
\usepackage{xcolor}
\usepackage{listings}
\usepackage{fancyhdr}

\pagestyle{fancy}
\fancyhf{}
\rhead{VSEPR-SIM v4.0.0}
\lhead{CleanGate \& Phase Models}
\rfoot{Page \thepage}

\lstset{
  language=C++,
  basicstyle=\ttfamily\small,
  keywordstyle=\color{blue!70!black},
  commentstyle=\color{green!50!black},
  stringstyle=\color{red!60!black},
  breaklines=true,
  frame=single,
  xleftmargin=1em,
  numbers=left,
  numberstyle=\tiny\color{gray}
}

\title{%
  \textbf{CleanGate Adaptive Verification} \\[0.3em]
  \textbf{\& Three-Tier Phase Models} \\[0.5em]
  \large VSEPR-SIM Methodology Document \\
  \large Section: Validation \& Thermodynamic Infrastructure
}
\author{VSEPR-SIM Project}
\date{April 2026 --- v4.0.0}

\begin{document}
\maketitle
\tableofcontents
\newpage

% ============================================================================
\section{Introduction}
% ============================================================================

This document describes two interconnected subsystems added to the VSEPR-SIM
atomistic platform in version 4.0.0:

\begin{enumerate}
  \item \textbf{CleanGate}: A lightweight adaptive verification gate that
        replaces rigid exhaustive validation with stratified semi-random
        sampling and risk-driven escalation.
  \item \textbf{Three-Tier Phase Models}: A selectable hierarchy of
        melt/gas behavior estimators at different fidelity levels.
\end{enumerate}

Both systems follow the anti-black-box design philosophy: every intermediate
score, every sampled module, and every escalation decision is explicitly
recorded and traceable.

% ============================================================================
\section{CleanGate: Adaptive Verification}
% ============================================================================

\subsection{Motivation}

Not every molecular or material case requires the same depth of validation.
A simple, stable, well-classified system should pass through quickly.
An exotic, thermally stressed, or poorly converging system should receive
deeper scrutiny.

CleanGate replaces the ``250-check religious ritual'' with:
\begin{itemize}
  \item Stratified semi-random sampling of gate modules across tiers
  \item Risk scoring via interpretable variables
  \item Escalation only when justified by measured irregularity
\end{itemize}

\subsection{Gate Tier Registry}

Modules are organized into 51 gates across five tiers:

\begin{table}[h]
\centering
\begin{tabular}{llll}
\toprule
\textbf{Tier} & \textbf{Range} & \textbf{Example Modules} & \textbf{Character} \\
\midrule
CORE   & 0--2   & identity, bounds, geometry       & Always run \\
LOW    & 3--10  & formatting, species, charge       & Basic sanity \\
MID    & 11--25 & geometry plausibility, bonding    & Chemistry checks \\
HIGH   & 26--50 & partial relaxation, stability     & Physics probes \\
TARGET & ---    & case-type specific                & Adaptive \\
\bottomrule
\end{tabular}
\caption{Gate tier classification with 51 registered modules.}
\end{table}

\subsection{Semi-Random Stratified Sampling}

The sampling strategy balances repeatability with coverage variety:

\begin{enumerate}
  \item Always include all 3 CORE modules (identity, bounds, geometry)
  \item Randomly choose 2 from LOW tier (T3--T10)
  \item Randomly choose 2 from MID tier (T11--T25)
  \item Randomly choose 2 from HIGH tier (T26--T50)
  \item Add 1 targeted module based on detected case type:
    \begin{itemize}
      \item Protein $\rightarrow$ thermal response tendency
      \item Organometallic $\rightarrow$ coordination number
      \item Exotic lattice $\rightarrow$ phonon stability
      \item Thermal stress $\rightarrow$ thermal expansion
      \item Complex organic $\rightarrow$ torsion range
    \end{itemize}
\end{enumerate}

Total: approximately 10 modules per evaluation (vs.\ 51 exhaustive).

\subsection{Weirdness Score}

The weirdness score $W$ measures how far a case departs from the normal
envelope of known, stable, classifiable systems:

\begin{equation}
  W = w_1 C_{\text{class}} + w_2 G_{\text{geom}} + w_3 E_{\text{energy}}
    + w_4 B_{\text{bond}} + w_5 R_{\text{ref}}
\end{equation}

\noindent where:
\begin{table}[h]
\centering
\begin{tabular}{clll}
\toprule
\textbf{Symbol} & \textbf{Component} & \textbf{Meaning} & \textbf{Default $w_i$} \\
\midrule
$C$ & class-uncertain    & Hard to classify                & 0.20 \\
$G$ & geom-irregular     & Geometry looks abnormal          & 0.25 \\
$E$ & energy-conflict    & Cheap models disagree            & 0.20 \\
$B$ & bond-ambiguity     & Bonding assignment unstable      & 0.15 \\
$R$ & reference-distance & Far from known material families & 0.20 \\
\bottomrule
\end{tabular}
\caption{Weirdness score components and default weights (sum = 1.0).}
\end{table}

\noindent Classification bands:
\begin{itemize}[nosep]
  \item $0.00$--$0.25$: ordinary
  \item $0.25$--$0.50$: unusual
  \item $0.50$--$0.75$: highly irregular
  \item $0.75$--$1.00$: exotic or poorly classifiable
\end{itemize}

\subsection{Temperature Score}

The temperature score $T_s$ is \emph{not} just temperature---it is thermal
severity relative to the material/system:

\begin{equation}
  T_s = a_1 \frac{T}{T_{\text{ref}}}
      + a_2 \frac{|dP/dT|}{S_{\text{expected}}}
      + a_3 \Phi_{\text{phase}}
      + a_4 \Phi_{\text{noise}}
\end{equation}

\noindent where:
\begin{table}[h]
\centering
\begin{tabular}{clll}
\toprule
\textbf{Symbol} & \textbf{Component} & \textbf{Meaning} & \textbf{Default $a_i$} \\
\midrule
$T/T_{\text{ref}}$ & temperature ratio & Normalized severity & 0.30 \\
$|dP/dT|/S$        & thermal gradient  & Response sensitivity & 0.25 \\
$\Phi_{\text{phase}}$ & phase risk     & Transition proximity & 0.25 \\
$\Phi_{\text{noise}}$ & thermal noise  & Perturbation instability & 0.20 \\
\bottomrule
\end{tabular}
\caption{Temperature score components and default weights (sum = 1.0).}
\end{table}

\noindent Classification bands:
\begin{itemize}[nosep]
  \item $0.00$--$0.30$: thermally mild
  \item $0.30$--$0.60$: notable thermal influence
  \item $0.60$--$0.80$: high thermal stress
  \item $0.80$--$1.00$: likely phase-risk or decomposition-risk regime
\end{itemize}

\subsection{Convergence Confidence}

Estimated from minimization quality signals:
\begin{itemize}[nosep]
  \item $> 0.80$: strong
  \item $0.55$--$0.80$: acceptable
  \item $0.35$--$0.55$: marginal
  \item $< 0.35$: suspicious
\end{itemize}

Inputs: maximum residual force, RMS force, step count, convergence flag.
Excessive step counts (>5000) incur a penalty indicating a difficult
energy landscape.

\subsection{Emergence Classification}

Three states separate emergence detection from convergence proof:

\begin{enumerate}[nosep]
  \item \textbf{NONE}: No meaningful emergence detected
  \item \textbf{VISIBLE\_WEAK}: Emergence visible but weakly verified (low convergence)
  \item \textbf{VISIBLE\_STRONG}: Emergence visible and strongly verified
\end{enumerate}

This prevents the classic failure mode: a system looks interesting,
everyone gets excited, and six hours later it turns out the solver was
hallucinating with confidence.

\subsection{Escalation Logic}

\begin{enumerate}
  \item If all prereqs pass and risk is low $\rightarrow$ continue normally
  \item If $W > 0.70$ or $T_s > 0.75$ $\rightarrow$ enable verbose tracking
  \item If convergence $< 0.45$, or exotic/protein/complex-organic case
        $\rightarrow$ require deep gate modules
  \item If deep gate required and convergence $< 0.30$ $\rightarrow$
        mark result as provisional
\end{enumerate}

\subsection{Reporting}

Two report formats:

\textbf{Internal} (dense, numeric, machine-readable):
\begin{lstlisting}
clean_gate.pass_basic              = true
clean_gate.weirdness_score         = 0.82
clean_gate.temperature_score       = 0.67
clean_gate.convergence_confidence  = 0.41
clean_gate.prereq_coverage         = 0.58
clean_gate.escalate_verbose        = true
clean_gate.require_deep_gate       = true
\end{lstlisting}

\textbf{External} (human-readable summary):
\begin{lstlisting}
CleanGate Summary:
  - Basic prerequisites: PASSED
  - Material/system irregularity: exotic (0.82)
  - Thermal severity: high thermal stress (0.67)
  - Convergence confidence: marginal (0.41)
  - Deep verification is recommended.
\end{lstlisting}

% ============================================================================
\section{Three-Tier Phase Models}
% ============================================================================

\subsection{Overview}

Three selectable modes provide melt and gas behavior estimation at
different fidelity/cost tradeoffs:

\begin{table}[h]
\centering
\begin{tabular}{llll}
\toprule
\textbf{Mode} & \textbf{Name} & \textbf{Melt Method} & \textbf{Gas Method} \\
\midrule
A & Threshold-Enthalpy      & Enthalpy threshold        & Ideal + $Z$-correction \\
B & Smoothed Phase-Fraction & Sigmoid/piecewise $f_\ell$ & Cubic EOS (VdW/RK) \\
C & Stability-Competition   & Gibbs $G_s$ vs $G_\ell$   & Fugacity (RK-based) \\
\bottomrule
\end{tabular}
\caption{Phase model modes and their methods.}
\end{table}

\subsection{Mode A: Threshold-Enthalpy (Fast Screening)}

\subsubsection{Melting}

Accumulated enthalpy:
\begin{equation}
  H(T) \approx \int_{T_0}^{T} c_p(T')\,dT' \approx c_{p,s} \cdot T
\end{equation}

Melt onset when $H(T) \geq H_{\text{melt,onset}}$; full melt when
$H(T) \geq H_{\text{melt,onset}} + \Delta H_f$.

\subsubsection{Gas Behavior}

Ideal gas with compressibility correction:
\begin{equation}
  PV = ZnRT, \qquad Z_{\text{est}} = 1 + \alpha\frac{P}{T} + \beta\rho
\end{equation}

\textbf{Best use}: very fast prechecks, rough material ranking,
large-scale parameter sweeps.

\textbf{Weakness}: melting treated too sharply; poor near phase boundaries.

\subsection{Mode B: Smoothed Phase-Fraction (Practical)}

\subsubsection{Melting}

Sigmoid liquid fraction:
\begin{equation}
  f_\ell(T) = \frac{1}{1 + \exp[-k(T - T_m)]}
\end{equation}

or piecewise linear:
\begin{equation}
  f_\ell(T) = \begin{cases}
    0 & T < T_s \\
    \frac{T - T_s}{T_\ell - T_s} & T_s \leq T \leq T_\ell \\
    1 & T > T_\ell
  \end{cases}
\end{equation}

Effective heat capacity with latent heat contribution:
\begin{equation}
  c_{p,\text{eff}} = (1 - f_\ell)\,c_{p,s} + f_\ell\,c_{p,\ell}
                    + \Delta H_f \frac{df_\ell}{dT}
\end{equation}

Effective density:
\begin{equation}
  \rho_{\text{eff}} = (1 - f_\ell)\,\rho_s + f_\ell\,\rho_\ell
\end{equation}

\subsubsection{Gas Behavior}

Cubic equations of state. Van der Waals:
\begin{equation}
  \left(P + \frac{a\,n^2}{V^2}\right)(V - nb) = nRT
\end{equation}

Redlich-Kwong:
\begin{equation}
  P = \frac{RT}{V_m - b} - \frac{a}{\sqrt{T}\,V_m(V_m + b)}
\end{equation}

Newton iteration solvers with deterministic convergence bounds.

\textbf{Best use}: transition-region modeling, materials that soften
before full melt, non-ideal gases.

\textbf{Weakness}: still empirical; needs fitted parameters.

\subsection{Mode C: Stability-Competition (Serious)}

\subsubsection{Melting}

Gibbs free energy for each phase:
\begin{equation}
  G_\phi(T, P) = H_{\phi,\text{ref}} + c_{p,\phi}(T - T_{\text{ref}})
               - T\left[S_{\phi,\text{ref}} + c_{p,\phi}\ln\frac{T}{T_{\text{ref}}}\right]
               + PV_\phi
\end{equation}

Melt-driving potential:
\begin{equation}
  \Delta G_{\text{melt}} = G_\ell - G_s
\end{equation}

Interpretation: $\Delta G > 0$ means solid favored; $\Delta G < 0$ means
liquid favored.

Probabilistic liquid fraction via Boltzmann softening:
\begin{equation}
  f_\ell = \frac{1}{1 + \exp(\Delta G_{\text{melt}} / \gamma)}
\end{equation}

\subsubsection{Gas Behavior}

Fugacity-based treatment using Redlich-Kwong EOS:
\begin{equation}
  \mu = \mu^\circ(T) + RT\ln f, \qquad f = \phi P
\end{equation}

Fugacity coefficient:
\begin{equation}
  \ln\phi = (Z-1) - \ln(Z - B') - \frac{A'}{B'}\ln\left(1 + \frac{B'}{Z}\right)
\end{equation}

where $A' = aP/(R^2 T^{5/2})$, $B' = bP/(RT)$.

\textbf{Best use}: serious phase modeling, exotic materials,
publication-grade results.

\textbf{Weakness}: parameter hungry; computationally heavier.

\subsection{Output Variables}

All three modes produce a unified \texttt{PhaseReport} with:

\begin{table}[h]
\centering
\begin{tabular}{lll}
\toprule
\textbf{Variable} & \textbf{Range} & \textbf{Meaning} \\
\midrule
\texttt{melt\_score}            & $[0,1]$ & Liquid fraction estimate \\
\texttt{gas\_nonideality\_score} & $|Z-1|$ & Departure from ideal gas \\
\texttt{phase\_risk\_score}      & $[0,1]$ & Proximity to phase transition \\
\texttt{temperature\_severity}   & $[0,1]$ & Normalized thermal stress \\
\bottomrule
\end{tabular}
\caption{Unified output variables across all phase model modes.}
\end{table}

% ============================================================================
\section{Integration with CleanGate}
% ============================================================================

The phase model feeds directly into CleanGate's temperature score:

\begin{itemize}
  \item \texttt{phase\_risk\_score} $\rightarrow$ $\Phi_{\text{phase}}$ in $T_s$
  \item \texttt{temperature\_severity} provides the $T/T_{\text{ref}}$ ratio
  \item \texttt{gas\_nonideality\_score} contributes to weirdness via
        energy conflict signals when multiple EOS modes disagree
\end{itemize}

The \texttt{melt\_score} also informs CleanGate's case type detection:
systems with $f_\ell \in (0.01, 0.99)$ are automatically tagged as
thermal stress cases for targeted module selection.

% ============================================================================
\section{File Manifest}
% ============================================================================

\begin{table}[h]
\centering
\begin{tabular}{lll}
\toprule
\textbf{File} & \textbf{Lines} & \textbf{Purpose} \\
\midrule
\texttt{atomistic/validation/clean\_gate.hpp}  & $\sim$600 & CleanGate system \\
\texttt{atomistic/core/phase\_model.hpp}       & $\sim$540 & Three-tier phase models \\
\texttt{docs/section\_clean\_gate\_phase\_models.tex} & this file & Methodology document \\
\bottomrule
\end{tabular}
\caption{Files created for CleanGate and phase model infrastructure.}
\end{table}

% ============================================================================
\section{Design Principles}
% ============================================================================

\begin{enumerate}
  \item \textbf{Deterministic over theatrical}: All scores are reproducible
        given the same seed and input data.
  \item \textbf{Anti-black-box}: Every weirdness component, every sampled
        module, every Gibbs energy intermediate is exposed and traceable.
  \item \textbf{Monotonic and interpretable}: Scores increase monotonically
        with increasing risk/severity. No inversion tricks.
  \item \textbf{Composable with existing infrastructure}: CleanGate follows
        the \texttt{GateResult}/\texttt{ValidationReport} pattern from
        \texttt{validation\_gates.hpp}. Phase models compose with
        \texttt{gas2\_eos.hpp} EOS types and \texttt{atomistic::thermo}
        constants.
  \item \textbf{No ``meso'' terminology}: All references use ``atomistic''
        per permanent project rule.
\end{enumerate}

\end{document}
